\section{Scale selection before trajectory evaluation}
\label{sec:selection_diagnostic}
We use $d\in\{1.6,1.8,2,2.2,2.4\}$ and
$h\in\{1.2,1.5,1.78,2.1,2.4\}$, with $a=1$, $r=1.29$,
$\alpha=.001$, $\tau=\nu=\alpha^2$, $v_L=\tau^2+\alpha^2$,
and $\lambda=1+\alpha^2$. The candidate grid is prespecified
relative to this evaluation. We use the S4 tolerances on $E$
and an allowance of $10^{-4}$ for the selection inequalities.
The trajectory simulations use the original RED update in double
precision.

\paragraph{Selection.}
We take the largest $V_-$ in $\{.08,.09,\dots,.30\}$ passing
the low-scale bound of S5 and the smallest $V_+$ in
$\{1.5,1.6,\dots,8\}$ passing its high-scale bound. Both bounds
cover the full restricted ranges $[v_L,V_-]\cup[V_+,\infty)$.
For the channel we take the first passing width-$.02$ band
in increasing order of its center in $\{.40,.42,\ldots,1.10\}$.
We check class-1 progress, transverse invariance, overshoot,
and invariance of the prepared strip through call 25.
The progress comparison starts from $\alpha-\alpha/4$.
Classes 0 and 2 start in their terminal rectangles and use only
$v_L$ under the observation-based preset. We take
$G_2=[-1.1,.01]\times[d-.27,d+.27]$ and use the main-text
$G_0,G_1$. We check the classifier, initialization, positive
template gaps at all corners, and every terminal face inequality.
The bound calculations use 45-digit outward-rounded interval
arithmetic. They hold on $E$ for each certified grid model.
The cases $(d,h)=(1.6,2.1)$ and $(1.6,2.4)$ do not pass a
candidate band under these tests.

\paragraph{Budget-based selection.}
The first-pass rule selected $[.39,.41]$ at all 23 passing
models. After that diagnostic we froze a second rule to
study model-dependent selection. It extends the candidate
centers down to $.24$, and among passing bands minimizes the
first call at which the analytic lower comparison reaches
$.53$, breaking ties by the smaller center. This uses the
model inputs and the 25-call cap. The actual experimental
prefix remains 25 calls for all policies. The new rule
certifies the same 23 settings, with selected centers
between $.36$ and $.60$ and upper bounds of 14 to 24 calls.
It is evaluated on a new noise stream. We execute the
shorter certified prefixes in the follow-up below.

\paragraph{Frozen evaluation.}
After writing each plan, we draw 32 independent pairs of
within-component and measurement noise per class and model,
using seed 20261112 for the first-pass rule and seed
20261113 for the budget-based rule. All 2208 observations
in each evaluation lie in $E$. All policies share their
observations, backprojection, step, terminal variance and
106-call budget. We use 25 prefix calls with alternating
restricted cutoffs, 25 geometric variances from $V_+$ to
$V_-$, or a constant platform at the selected center.
The preset assigns the constant platform only when the
observation classifier predicts class 1. Class labels enter
evaluation and prior weighting. The original terminal map
is then used for at least 81 calls. We retain first entry
times, class counts at the common switch, and class and model
squared errors. Model averages describe the certified grid
subset. The other two models remain uncertified.

Both designs yield the same class-entry counts: restricted
switching has zero passing models, both shared geometric and
shared-platform prefixes pass class 1 on 23 models but not all
classes, and the observation preset passes all classes on all 23.
The first-pass mean errors in this order are
$.32154626,.24926068,.24926068$ and $1.5293\cdot10^{-12}$.
Every class-1 draw under the budget-based rule reaches $G_1$
within its certified bound. These are allocation ablations;
the current main-text table instead reports the independent,
budget-matched controls of S29.

The prespecified failure examples both use $(d,h)=(2,1.78)$.
They hold $V_+=2.2$ for 12 calls and then use $V_-=.18$,
with or without one call at the selected center at the start
of the low stage. Both produce zero class-1 terminal entries
in 106 calls. For the budget-based band $[.43,.45]$, one band
call leaves class 1 at $s<.194$ at the switch. The shared
geometric grid reaches the correct class-1 rectangle on this
fixture at call 19 without using a variance in $[.43,.45]$.
The experiment checks a sufficient prescription, and nearby
successful schedules can use other intermediate variances.

\subsection{Independent platform tuning and realized savings}
We freeze an additional plan on the same 23 certified grid
models. The shared and observation-gated platforms have
identical model inputs. Each searches
$v\in\{.20,.21,\ldots,1.10\}$ and prefix lengths $m=1,\ldots,25$
independently, giving $91\times25=2275$ candidates per
model and policy. For each candidate, run RED from all
three noiseless backprojections $B\mu_z$ with 106 total calls.
The gated policy uses the same classifier as the certified
preset. A candidate passes when the maximum final squared
distance to the corresponding template is at most $10^{-8}$.
Select the shortest passing prefix and the lower median
of its passing variances. Every search has a passing
candidate. The shared choices have $v\in[.41,.85]$ and
the gated choices $v\in[.47,.81]$. Both choose 11--14 calls.

On the noiseless grid the shared search passes 9--25
of the 2275 pairs per model, versus 428--1142 for gating.
At fixed length 25, the counts are 0--2 of 91 for sharing
and 48--86 of 91 for gating. These are discrete measures
of sensitivity to the tuning choices on the model proxies.
Only the frozen choices are carried to the noisy evaluation,
apart from the explicit nominal comparison at $v=.25$.

With seed 20261114 we draw new independent pairs of component
and measurement noise, 64 per class and model, except for the
nominal $(2,1.78)$ model where we draw 512 per class.
All 5760 observations lie in $E$. The certified preset
uses its selected center and analytic cap $m_{\rm cert}$.
All main comparisons use the same 106 calls, observations,
initialization, exact RED formula and terminal variance.
The shorter-budget variant stops at $m_{\rm cert}+81$ on the
same observations. The class labels are used for evaluation
and prior weighting; the sample-specific gate sees only $b$.

\begin{table}[t]
\centering\small
\setlength{\tabcolsep}{2.8pt}
\begin{tabular}{@{}lrrrr@{}}
\toprule
Prefix policy & $G_1$ & All $G_z$ & Mode & Error\\
\midrule
Shared, same center & 23 & 0 & 0 & $.249$\\
Preset, 25-call cap & 23 & 23 & 23 & $1.47\cdot10^{-12}$\\
Preset, certified cap & 23 & 23 & 23 & $1.47\cdot10^{-12}$\\
Shared, tuned & 0 & 0 & 23 & $1.47\cdot10^{-12}$\\
Gated, tuned & 0 & 0 & 23 & $1.47\cdot10^{-12}$\\
Preset, shorter budget & 23 & 23 & 23 & $1.47\cdot10^{-12}$\\
\bottomrule
\end{tabular}
\caption{Frozen follow-up on 23 models. Columns $G_1$ and All
$G_z$ count models whose indicated draws all lie in the named
rectangles after the prefix window. Mode counts models with every
draw's final nearest template correct. Error is prior-weighted
squared error averaged across models.
All rows use 106 total calls except the last (95--105 calls).}
\label{tab:tuned_platforms}
\end{table}

The sufficient rectangles are stricter than the final nearest
template test. On the nominal model, the shared 25-call platform
at total variance $.25$ switches with horizontal coordinates
near $.4961$ and $.5018$ for classes 0 and 1, outside $G_0$ and
$G_1$. It still yields $1536/1536$ correct nearest templates and
prior-weighted squared error $1.5020\cdot10^{-12}$, while the
shared $.44$ prefix yields $1024/1536$ and $.24521$. The tuned
shared choice $(v,m)=(.61,12)$ also yields $1536/1536$ and
$1.5020\cdot10^{-12}$. The shared outcomes are numerical
checks on the selected grid, without an interval guarantee.

For the certified policy, the mean prefix cap drops from 25
to 17.22. Averaging the actual intermediate calls first with
the class priors and then over models gives 8.05 and 5.54 calls
before and after this change. The shorter run saves a mean of
7.78 total calls at the observed error in Table~\ref{tab:tuned_platforms}.

\paragraph{Reproduction.}
The two phases of
\path{experiments/code/scale_selection_diagnostic.py}
write a plan and evaluate it from a separate command.
The frozen plans and per-model outputs are in
\path{experiments/results/scale_selection} and
\path{experiments/results/scale_selection_budget}.
The follow-up uses \path{experiments/code/scale_selection_followup.py}
with \texttt{plan} and then \texttt{evaluate}, both with
\path{--out experiments/results/scale_selection_followup}
when run in a new artifact tree without the reference outputs.
The first command reads the frozen budget-based plan by default.
The second reports the policy summaries in
Table~\ref{tab:tuned_platforms} and writes the per-model output.
